\secciongrande{Programación de scripts dónde y cómo colocar los scripts}

\resumen{
  Los scripts de javascript se insertan en un documento html utilizando la etiqueta \texttt{<script>}. Pueden ubicarse tanto en el encabezado (entre \texttt{<head>} y \texttt{</head>}) como en el cuerpo (entre \texttt{<body>} y \texttt{</body>}). La etiqueta de apertura debe incluir el atributo \texttt{type="texto/javascript"} para indicar el formato de texto javascript.
}

\definicion{Etiqueta script}{Etiqueta contenedora html que designa el inicio y fin de un bloque de código javascript, comenzando con \texttt{<script>} y terminando con \texttt{</script>}.}

\subtituloA{Ubicaciones válidas en html}

\lista{
  \textbf{Encabezado:} Colocación entre \texttt{<head>} y \texttt{</head>}, destinada a metadatos o scripts que deben cargarse antes del contenido visible. |
  \textbf{Cuerpo:} Colocación entre \texttt{<body>} y \texttt{</body>}, ideal para scripts que interactúan con elementos visibles de la página.
}

\nota{Para el examen}{
¿Cuál es la función del atributo \texttt{type="texto/javascript"} en la etiqueta \texttt{<script>}? Indica al navegador que el contenido es código javascript en formato de texto, no html.
}
